% IT-Security ToDo-Liste
% Aspekte für den Einsatz in einem datenschutzsensiblen Umfeld
% Fokus: GUI_4 (BMS-GUI), PostgreSQL, pgAdmin4, Docker-Umgebung

\textbf{Offene IT-Security- und Datenschutz-Aspekte}

Einführung eines rollenbasierten Benutzer- und Rechtekonzepts für die GUI (Authentifizierung, Autorisierung, Session-Handling).
Absicherung der GUI-Kommunikation über HTTPS (TLS), inklusive gültiger Zertifikate und HSTS-Konfiguration.
Trennung von Entwicklungs-, Test- und Produktionsumgebung mit jeweils eigenen Datenbanken und Zugangsdaten.
Vertrauliche Speicherung aller Zugangsdaten (DB-User, Passwörter, API-Keys) in Config/Secrets-Management statt im Quellcode oder in Klartext-Dateien.
Einschränkung der PostgreSQL-Benutzerrechte (Least-Privilege-Prinzip: getrennte Rollen für Lesen, Schreiben, Administration).
Härtung der PostgreSQL-Konfiguration (z.B. pg\_hba.conf, sslmode, Logging, Passwort-Policy, Brute-Force-Schutz über Fail2Ban/Firewall-Regeln).
Absicherung und Härtung von pgAdmin4 (starkes Admin-Passwort, IP-Restriktionen, keine Nutzung als allgemeines User-Frontend).
Verschlüsselung sensibler Daten in der Datenbank bzw. Verschlüsselung der Datenträger (Full-Disk-Encryption oder Tablespace-Verschlüsselung).
Konzeption und Umsetzung eines Berechtigungskonzepts für personenbezogene Daten (Wer darf welche Kandidaten- und Jobdaten sehen/verarbeiten?).
Ergänzung von Logging- und Audit-Funktionen (Wer hat wann welche Daten gelesen, geändert oder exportiert?).
Implementierung eines strukturierten Lösch- und Sperrkonzepts (inkl. Löschfristen und Dokumentation gemäß DSGVO).
Prüfung von Datenminimierung und Zweckbindung (nur notwendige personenbezogene Felder speichern, keine überflüssigen Freitexte).
Umgang mit Backups: verschlüsselte Backups, sichere Aufbewahrung, definierte Aufbewahrungsdauer und geregelte Wiederherstellungstests.
Härtung der Docker-Container (kein Root-User im Container, minimale Basis-Images, regelmäßige Sicherheitsupdates und Image-Scanning).
Absicherung des Docker-Netzwerks (Netzwerksegmentierung, keine unnötig offenen Ports, Zugriff nur aus definierten Netzen).
Monitoring und Alerting für Datenbank, Container und Host-System (Ressourcen, Fehlversuche, ungewöhnliche Zugriffsmuster).
Input-Validierung und -Sanitization an allen Benutzereingaben in der GUI (Schutz vor SQL-Injection, XSS, Command-Injection).
Schutz der LLM-/Ollama-Schnittstellen vor unautorisiertem Zugriff (nur intern erreichbar, Rate-Limiting, Authentifizierung falls extern exponiert).
Prüfung und Dokumentation der eingesetzten Drittkomponenten (Lizenzen, Security-Updates, bekannte Schwachstellen / CVEs).
Erstellung einer Datenschutz-Folgenabschätzung (DSFA) und eines Verzeichnisses der Verarbeitungstätigkeiten für den Produktivbetrieb.
Schulung der Administratoren und Fachanwender zu sicheren Betriebs- und Nutzungsprozessen (Passwort-Policy, Export/Weitergabe von Daten).
Regelmäßige Durchführung von Penetrationstests bzw. Sicherheitsreviews der Gesamtumgebung (GUI, Datenbank, Docker-Host, Netzwerk).
Notfall- und Incident-Response-Konzept (Erkennung, Meldung, Analyse und Kommunikation von Sicherheitsvorfällen und Datenpannen).
Technische Maßnahmen zur Anonymisierung/Pseudonymisierung von Daten für Test- und Analysezwecke (keine echten Kandidaten-Daten in Testsystemen).
Überprüfung und ggf. Ergänzung von Log- und Export-Funktionen hinsichtlich DSGVO-Konformität (z.B. wer darf Exportdateien erzeugen und wohin speichern?).
